\section{Conclusion}
\label{sec:conclusion}

We have presented graph learning for dynamic wireless topologies through a single
unifying lens: classical message passing, graph diffusion, and quantum walks are one
propagation operator $\Prop(\Lap)$ with different spreading physics, diffusive versus
ballistic. From this lens we derived a param-matched quantum-inspired GNN layer that
runs today, and we used a fully reproducible LEO case study to teach not only how to
build it but how to judge it. The case study's honest verdict is itself the lesson. Global propagation beats local by
a wide margin at every scale; and among global operators the quantum-inspired and
classical kernels trade places as the graph grows, the quantum-inspired walk losing at
$66$ satellites, tying at $264$, and becoming the best operator at $1584$, the
diameter-driven trend its ballistic spreading predicts. The value of a quantum mechanism
must be measured, not assumed, and the measurement must survive param-matching, multiple
seeds, scale-free metrics, a sweep of the obvious design choices, and a test at the scale
where the mechanism is supposed to help. We hope a reader leaves able to place any new quantum graph model
in the operator landscape and to evaluate it without fooling themselves.

The practical takeaways are short. (i) Treat GCN, diffusion, and quantum walks as one
propagation operator $\Prop(\Lap)$ and pick the spreading physics the task needs. (ii) On
a long-range task over a large dynamic graph, a global operator is mandatory and a fixed
$K$-hop GCN is the wrong tool. (iii) Among global operators, classical heat is a strong,
cheap default at small-to-medium scale; reach for a quantum-inspired walk on large,
long-range graphs, where our largest-scale experiment shows it becomes the best operator,
and confirm any gain with a multi-seed comparison. (iv) True
quantum hardware is a research path, not a deployment choice today: the ballistic
mechanism is real but fragile to NISQ noise. (v) Above all, evaluate honestly,
param-matched operators, multiple seeds, scale-free metrics, honest denominators, and a
design sweep before any verdict.

\section*{Data and Code Availability}
All experiments, the LEO topology generator, the operators, and the figure-generation
scripts are available at \url{https://github.com/haodpsut/qi-gnn-dynamic-wireless}.
Every figure in this tutorial is regenerated from seeded runs by the scripts in that
repository.
